\subsection{Hilbert schemes and the exact positive theta input}

\paragraph{P-F53 (a geometric coefficient model for the identified input).}
Let $S$ be a complex projective K3 surface, and let $S^{[m]}$
be its Hilbert scheme of length-$m$ subschemes, with
$S^{[0]}=\operatorname{Spec}\mathbb C$.
Retain the positive-definite lattice $K=E_6$, its discriminant
cosets $\mu\in K^\vee/K$, and the quadratic form
$Q(x)=(x,x)/2$ of P-F46 and P-F51. Write the actual
Borcherds input from P-F46 as
\[
 g_\mu(\tau)=\frac{\Theta_{K+\mu}(\tau)}{\delta(\tau)}
        =\sum_a c_\mu(a)q^a,
 \qquad \delta(\tau)=q\prod_{r\ge1}(1-q^r)^{24}.
                                                        \tag{P-F53a}
\]
For each rational exponent $a$, define the following finite disjoint union:
\[
 X_{\mu,a}=
 \coprod_{\substack{x\in K+\mu\\
                  m=a+1-Q(x)\in\mathbb Z_{\ge0}}} S^{[m]}.
                                                        \tag{P-F53b}
\]
An empty disjoint union is allowed. Then $X_{\mu,a}$ is a
smooth projective complex scheme (possibly disconnected and
of different dimensions on its components), its odd rational
cohomology vanishes, and
\[
          c_\mu(a)=\dim_{\mathbb Q}H^*(X_{\mu,a},\mathbb Q).
                                                        \tag{P-F53c}
\]
Each $H^j(X_{\mu,a},\mathbb Q)$ is a pure Hodge structure of
weight $j$. The statement concerns purity separately in each
cohomological degree, and does not identify all the cohomology
with one pure weight or assert that it is Tate.

\paragraph{The cited Hilbert-scheme formula and the proof.}
The smoothness and projectivity of $S^{[m]}$ are the
surface Hilbert-scheme theorem recalled in G\"ottsche,
\emph{Hilbert Schemes of Points on Surfaces}, Section~1,
printed p.~484. His Betti-number formula is equation~(2.1),
printed p.~485. In his sign convention the Poincar\'e
polynomial is $p(Y,z)=\sum_j(-1)^jb_j(Y)z^j$.
For a K3 surface the Betti numbers are $(1,0,22,0,1)$,
and the formula specializes to
\[
 \sum_{m\ge0}p(S^{[m]},z)q^m
 =\prod_{r\ge1}
   \frac1{(1-z^{2r-2}q^r)
           (1-z^{2r}q^r)^{22}
           (1-z^{2r+2}q^r)}.
                                                        \tag{P-F53d}
\]
All exponents of $z$ on the right are even. Comparing
coefficients of $q^m z^{2j+1}$ gives
$b_{2j+1}(S^{[m]})=0$.
At $z=1$ the signed polynomial therefore equals the
total cohomology dimension, and
\[
 \sum_{m\ge0}\dim H^*(S^{[m]},\mathbb Q)q^m
                   =\prod_{r\ge1}(1-q^r)^{-24}.
                                                        \tag{P-F53e}
\]
The same Euler-number specialization is explicitly
Corollary~2.4, printed p.~62, of G\"ottsche,
\emph{Invariants of Moduli Spaces and Modular Forms};
the K3 value $e(S)=24$ is stated there on p.~66.

For any fixed $a$, a summand in (P-F53b) requires
$Q(x)\le a+1$. Positive definiteness of $K$ makes that
set of coset vectors finite. For each such vector the
integer $m$ is unique. Hence (P-F53b) is a finite
disjoint union of smooth projective schemes. Finite
disjoint unions preserve smoothness and projectivity;
their rational cohomology is the direct sum of the
component cohomologies. This proves odd-degree
vanishing and the asserted degreewise purity, using
the usual pure Hodge structure on smooth projective
cohomology.

Finally, multiplying (P-F53e) by
$q^{-1}\sum_{x\in K+\mu}q^{Q(x)}$ and taking the
coefficient of $q^a$ gives precisely the direct sum
in (P-F53b), so proves (P-F53c). All coefficient sums
are finite; there is no interchange of conditionally
convergent series. This completes the proof of the
coefficient model.\hfill$\square$

\paragraph{The completed lambda-ring identity.}
This construction also gives an exact, limited lambda-ring
dictionary for the inverse-discriminant operation. Let
$V=H^*(S,\mathbb Q)$ after forgetting its cohomological
grading, so $\dim V=24$. In the completed Grothendieck
ring of finite-dimensional rational vector spaces, put
$\sigma_t(V)=\sum_{j\ge0}[\operatorname{Sym}^jV]t^j$.
Then
\[
 \sum_{m\ge0}[H^*(S^{[m]},\mathbb Q)]q^m
   =\prod_{r\ge1}\sigma_{q^r}(V)
   =\operatorname{Sym}\!\left(V\otimes q\mathbb Q[q]\right)
                                                        \tag{P-F53f}
\]
as classes graded by the exponent of $q$. The last term
denotes the graded class of the symmetric algebra, with
one copy of $V$ in each positive degree. To prove the
identity, the generating function of the dimensions
of symmetric powers of a 24-dimensional vector space
is $(1-t)^{-24}$: a monomial basis counts weak
compositions into 24 parts. Its product over $r$
is (P-F53e). Dimension identifies the Grothendieck
ring of finite-dimensional rational vector spaces
with $\mathbb Z$, so equality of these dimensions
proves the stated class equality in every grade.
Every grade receives only finitely many contributions.
Thus no infinite-dimensional object is being inserted
as one element of the uncompleted ring.

Equation (P-F53f) uses the symmetric operations
$\sigma_t=\lambda_{-t}^{-1}$ in the sense explained in
the Grothendieck plethysm bridge. It is an equality
after forgetting the cohomological grading. This
elementary proof does not assert a canonical
cohomological-ring isomorphism or compatibility with
the nonstandard algebra's multiplication. The much
stronger geometric Heisenberg construction is not
needed for this statement.

\paragraph{The exact principal coefficients and what has been identified.}
The previously checked counts for the $E_6$ cosets
make the nonpositive part completely geometric:
\[
 X_{0,-1}=\operatorname{Spec}\mathbb C,\qquad
 X_{0,0}=S\ \coprod\ 72\operatorname{Spec}\mathbb C,
 \qquad
 X_{\pm,-1/3}=27\operatorname{Spec}\mathbb C.
                                                        \tag{P-F53g}
\]
Indeed the zero coset has one vector of norm zero
and 72 vectors with $Q=1$, while each other coset
has 27 vectors of minimal $Q=2/3$, as proved in
P-F46. Their total dimensions are $1,96,27,27$.
In particular the weight $96/2=48$ of the identified
Borcherds output can be read as half the total
cohomology dimension of $X_{0,0}$.

This supplies smooth projective coefficient spaces
for the \emph{entire particular input} of P-F46,
and expresses its inverse-discriminant factor by
completed symmetric operations on an actual
surface cohomology. Together with P-F51, both
factors now have specified geometric sources.
It does not turn the regularized multiplicative
Borcherds lift into the supplied iterated definite
unitary lift, nor give a cohomological map from
these spaces onto the cubic multiplicity spaces.
The coefficient-growth obstruction of P-F46 still
rules out identifying the whole input with all
$T^{(5)}$ grades by a positive affine reindexing.
The pure Hodge pieces here also need not be Tate:
already $S\subset X_{0,0}$ has its nonzero
$(2,0)$ part. This example therefore adds an
actual Hilbert-scheme and lambda-operation source
to the remaining operation-comparison problem,
without claiming that purity alone solves it.

Primary sources for the cited formulae:
\url{https://arxiv.org/pdf/math/0304302},
Section~1 p.~484 and equation~(2.1) p.~485;
\url{https://rendiconti.dmi.units.it/volumi/41/06.pdf},
Corollary~2.4 p.~62 and the K3 specialization p.~66.
